% Avsnitt 18.10, ur lectures/L18/README.md.

\section{Sammanfattning}
\label{c18:sec:review}

\needspace{6\baselineskip}
\subsection*{Det här ska du kunna nu}
\begin{itemize}
  \item Förklara skillnaden mellan en puls och en klibbig nivå, och varför registerlagret måste
    konvertera mellan dem.
  \item Implementera \code{register_bank.vhd} enligt registerkartan: STATUS-bitarnas sättning och
    nollställning, \code{TX_SEND} som encykelspuls, och maskningen vid skrivning.
  \item Läsa ett SPI-tidsdiagram i läge 0 och säga vilken bit som ligger på MOSI vid varje flank.
  \item Förklara varför en SPI-slav i en FPGA översamplar SCK i stället för att använda den som
    klocka.
  \item Läsa och använda \code{spi_slave.vhd} utan att skriva den.
\end{itemize}

\needspace{6\baselineskip}
\subsection*{Frågor att testa dig själv med}
\begin{enumerate}
  \item Varför kan en drivrutin som pollar var femte mikrosekund inte se en puls på 20~ns?
  \item Vad sätter STATUS bit 0, och vad nollställer den?
  \item Varför låser felbiten på en \emph{stigande flank} av \code{error} och inte på nivån?
  \item En andra ram tas emot innan den första kvitterats. Vad händer, och varför är det ett medvetet
    val snarare än en bugg?
  \item STATUS bit 0 sätts av tre saker under drift, utöver reset. Vilka, och varför räcker det inte
    med \code{tx_done} ensamt?
  \item Vad gör \code{TX_ABORT}, och vad når det \emph{inte}?
  \item En skrivning av \code{0xFFFFFFFF} till \code{TX_DLC} läses tillbaka som \code{0x0000000F}.
    Var sker maskningen, och varför validerar banken inte att värdet är högst 8?
  \item Varför använder \code{spi_slave} inte SCK som klocka? Vad skulle gå fel om den gjorde det?
  \item Varför byter MISO värde på SCK:s fallande flank?
\end{enumerate}

\needspace{6\baselineskip}
\subsection*{Riktvärde}
\code{register_bank} är en avgränsad modul med en fullständig testbänk och bör gå att bli klar med
under passet och veckan efter. Den beror inte på SPI alls, så den kan skrivas parallellt med att
någon annan i gruppen läser in sig på transaktionsformatet. Lägg in den utdelade
\file{spi_slave.vhd} i repot oförändrad, och läs igenom \file{spi_reg_bridge_tb.vhd} inför nästa
kapitel: den är kontraktet för modulen du skriver då.

\needspace{6\baselineskip}
\subsection*{Blicken framåt}
\Chapref{c19:ch} bygger SPI-bryggan, toppnivån \code{can_spi_node}, och tar hela kedjan ut på riktig
hårdvara. Det är projektets sista pass.
